% TOP_PROBLEM: {"problemId": "problem.provisional-top500-omission-wave-004.018", "canonicalTitle": "Palis finite-attractor full-basin density conjecture", "releaseRank": 205, "primaryDomain": "probability_ergodic_dynamics", "primaryDomainLabel": "Probability, ergodic theory & dynamics", "displayStatus": "open", "releaseStatus": "open"}
% !TeX program = lualatex
% Definition and sources only; this file is not an LLM solution attempt.
\documentclass[11pt]{article}
\usepackage[margin=25mm]{geometry}
\usepackage{fontspec}
\setmainfont{Latin Modern Roman}
\usepackage{amsmath,amssymb,mathtools}
\usepackage{unicode-math}
\setmathfont{Latin Modern Math}
\usepackage{xurl,hyperref}
\hypersetup{colorlinks=true,urlcolor=blue}
\setcounter{secnumdepth}{0}
\setlength{\parindent}{0pt}
\setlength{\parskip}{5pt}
\setlength{\emergencystretch}{3em}
\begin{document}
\section*{\texorpdfstring{Palis finite-attractor full-basin density conjecture}{Palis finite-attractor full-basin density conjecture}}\label{205}
\subsection{Definitions and mathematical statement}
For a diffeomorphism $f$ of a compact manifold, let
$B(A)=\{x:\operatorname{dist}(f^n(x),A)\to0\}$; for a flow replace $n$ by continuous positive time. In the selected Palis convention an attractor is a compact invariant transitive set with positive-volume basin, not necessarily an attractor defined by a trapping neighborhood. For each finite $r\ge1$, the target is density in both $\operatorname{Diff}^r(M)$ and the space of $C^r$ vector fields of systems with finitely many such sets $A_1,\ldots,A_k$ satisfying
\[
 \operatorname{vol}\left(M\setminus\bigcup_{i=1}^k B(A_i)\right)=0.
\]
The number $k$ may vary with the system. Both density assertions are required; physical measures and stochastic stability are not included in this selected part of the global conjecture.
\par\smallskip\textit{Formulation and concept references:} \href{https://ems.press/content/serial-article-files/15975}{[S1]}.

\subsection{Short English statement}
Can every smooth dynamical system be approximated by one in which almost every initial point is attracted to one of finitely many transitive attractors?
\subsection{Sources}
\begin{itemize}
\item[S1]J. Palis, A global perspective for non-conservative dynamics, Ann. IHP AN 22 (2005), 485-507.
\url{https://ems.press/content/serial-article-files/15975}.
\textit{Formulation source.}
\end{itemize}
\end{document}
